\section{A marker-free resolution-and-component PRE-A member}

\textsc{PIVOT-FACTOR1} adds a second causal construction to the WRAPPER
portfolio.  Unlike the marked cutset compiler, it needs no supplied
certificate.  From source syntax it tries a Davis--Putnam pivot, generates all
nontautological resolvents, finds independent components, and dry-runs their
complete dense-CNF1 encodings.  Only after verifying local work, maximum-child
and aggregate-volume caps does it execute the children and AND their answers.
UNKNOWN makes zero SAT calls, but all failed discovery work is still charged.

On every accepted input the exact bounds are
\[
 T_{\rm local}=N^{6/5+o(1)},\qquad
 \max_i|Q_i|=N^{1/2+o(1)},\qquad
 \sum_i|Q_i|=N^{11/10+o(1)},
\]
with simultaneous space $N^{1/2+o(1)}$ including the live child.  Repeated
child encodings count with multiplicity.  A connected infinite family with a
positive hub is accepted: eliminating the hub exposes linearly many constant
size independent components.  Thus the member is nonvacuous and genuinely
marker-free.

It is not all-input.  The family
$\bigwedge_i(x\vee a_i)\wedge\bigwedge_j(\neg x\vee b_j)$ defeats every
one-pivot choice.  Connected regular parity CNFs also defeat the standard
unit/pure/non-growing-BVE extension: every pivot expands the exact byte
encoding and factoring yields one source-scale component.  Explicit complete
Davis--Putnam elimination on expander parity instances can require
exponential intermediate clauses, so allowing uncontrolled transient growth
does not meet the favorable ledger.  Compact XOR/Gaussian reasoning is a
separate typed mechanism and remains a legitimate next target.

This construction enlarges PRE-A coverage but supplies neither recursive
closure nor a favorable fallback on its UNKNOWN domain.  The exact unresolved
implication remains SAT-in-P yielding one uniform source-first favorable
phase on every input and descendant.  Conditional $K$ elimination is
preserved; $P$ versus $NP$ remains \textsc{Open}.

